\chapter{跨学科综合例题与面试压缩}

\begin{mapbox}
本章不是公式汇编，而是一组适合在纸上重新做一遍的完整例题。每题都从模型假设开始，写出决定性中间步骤，并在末尾检查量纲、极限与物理图像。建议第一次遮住解答独立完成，第二次把解答压缩成两分钟口述，第三次只在黑板上写主线。
\end{mapbox}

\section{理论力学：从几何约束到正则结构}

\begin{examplebox}[title=例 1：可滑动支点上的单摆]
质量为 $m$ 的小球由长为 $\ell$ 的轻杆连接到质量为 $M$、可沿水平面无摩擦滑动的支点。取支点坐标 $X$ 和摆角 $\theta$ 为广义坐标。小球位置为
\[
x=X+\ell\sin\theta,\qquad y=-\ell\cos\theta .
\]
逐项求导得到
\[
\dot x=\dot X+\ell\dot\theta\cos\theta,\qquad
\dot y=\ell\dot\theta\sin\theta .
\]
因此
\[
T=\frac12M\dot X^2+\frac12m(\dot x^2+\dot y^2)
=\frac12(M+m)\dot X^2+m\ell\dot X\dot\theta\cos\theta
+\frac12m\ell^2\dot\theta^2,
\]
\[
V=mg\ell(1-\cos\theta),\qquad L=T-V.
\]
$X$ 不显含于 $L$，故为循环坐标：
\[
P_X=\pdv{L}{\dot X}=(M+m)\dot X+m\ell\dot\theta\cos\theta=\text{常数}.
\]
这正是系统的水平总动量。对 $\theta$ 使用 Euler--Lagrange 方程，
\[
\pdv{L}{\dot\theta}=m\ell\dot X\cos\theta+m\ell^2\dot\theta,
\]
\[
\dv{}{t}\pdv{L}{\dot\theta}
=m\ell\ddot X\cos\theta-m\ell\dot X\dot\theta\sin\theta
+m\ell^2\ddot\theta ,
\]
\[
\pdv L\theta=-m\ell\dot X\dot\theta\sin\theta-mg\ell\sin\theta .
\]
相减后交叉项消去，得到
\[
\ell\ddot\theta+\ddot X\cos\theta+g\sin\theta=0.
\]
连同 $\dv{P_X}{t}=0$ 即为两个耦合运动方程。极限 $M\to\infty$ 给出固定支点单摆方程；这也是答案的必要检查。
\end{examplebox}

\begin{examplebox}[title=例 2：二自由度耦合振子的质量加权正规模]
两个质量均为 $m$ 的质点，各以弹簧 $k$ 连到墙，并以弹簧 $\kappa$ 相连。平衡位置附近
\[
T=\frac m2(\dot x_1^2+\dot x_2^2),\qquad
V=\frac k2(x_1^2+x_2^2)+\frac\kappa2(x_2-x_1)^2 .
\]
故
\[
M=m\begin{pmatrix}1&0\\0&1\end{pmatrix},\qquad
K=\begin{pmatrix}k+\kappa&-\kappa\\-\kappa&k+\kappa\end{pmatrix}.
\]
令 $\bm x=\bm a e^{i\omega t}$，运动方程化为
\[
(K-\omega^2M)\bm a=0.
\]
非零解要求
\[
0=\det(K-\omega^2M)
=(k+\kappa-m\omega^2)^2-\kappa^2.
\]
于是
\[
\omega_+^2=\frac{k}{m},\quad \bm a_+\propto(1,1);
\qquad
\omega_-^2=\frac{k+2\kappa}{m},\quad \bm a_-\propto(1,-1).
\]
同相模不拉伸中间弹簧，所以频率与 $\kappa$ 无关；反相模使中间弹簧伸长加倍，恢复力更强。把该思路推广到无穷晶格，就是声子正规模。
\end{examplebox}

\begin{examplebox}[title=例 3：由生成函数得到交换坐标和动量的正则变换]
取第二类生成函数
\[
F_2(q,P)=qP+\frac{\alpha}{2}q^2.
\]
定义
\[
p=\pdv{F_2}{q}=P+\alpha q,\qquad
Q=\pdv{F_2}{P}=q.
\]
所以 $Q=q,\ P=p-\alpha q$。直接检验基本 Poisson 括号：
\[
\{Q,P\}_{q,p}
=\pdv Qq\pdv Pp-\pdv Qp\pdv Pq
=1\cdot1-0\cdot(-\alpha)=1,
\]
且 $\{Q,Q\}=\{P,P\}=0$，因此变换保持辛结构。若旧 Hamiltonian 为
$H=p^2/(2m)+m\omega^2q^2/2$，新 Hamiltonian 为
\[
K(Q,P)=\frac{(P+\alpha Q)^2}{2m}+\frac12m\omega^2Q^2.
\]
生成函数不是任意换元的装饰；它保证正则一形式之差为全微分，从而保证 Hamilton 方程的形式不变。
\end{examplebox}

\section{电动力学：边界、传播与辐射}

\begin{examplebox}[title=例 4：接地导体球外点电荷的镜像解]
半径为 $a$ 的接地导体球外，在距球心 $b>a$ 处放电荷 $q$。设镜像电荷位于同一直线上，位置 $b'=a^2/b$，大小 $q'=-aq/b$。候选势为
\[
\phi(\bm r)=\frac1{4\pi\varepsilon_0}
\left(\frac q{\abs{\bm r-b\hat{\bm z}}}
+\frac{q'}{\abs{\bm r-b'\hat{\bm z}}}\right).
\]
在球面 $r=a$ 上，
\[
\abs{a\hat{\bm r}-b'\hat{\bm z}}^2
=a^2+\frac{a^4}{b^2}-2\frac{a^3}{b}\cos\theta
=\frac{a^2}{b^2}
\abs{a\hat{\bm r}-b\hat{\bm z}}^2.
\]
故第二个距离等于第一个距离的 $a/b$ 倍，且 $q'=-aq/b$，两项严格抵消，$\phi(a,\theta)=0$。候选势在球外除真实电荷处满足 Poisson 方程，并满足无穷远条件。由 Dirichlet 唯一性定理，它就是唯一物理解。

真实电荷受力等于镜像电荷给出的库仑力：
\[
F_z=-\frac1{4\pi\varepsilon_0}
\frac{a b q^2}{(b^2-a^2)^2}.
\]
负号表示吸引。不能把镜像电荷解释为导体内部真实存在的点电荷；它只是在求解区域外制造正确边界值。
\end{examplebox}

\begin{examplebox}[title=例 5：垂直入射时由边界条件推出 Fresnel 系数]
电磁波从介质 1 垂直入射到介质 2。设两介质无自由面电荷和面电流，波阻抗
$Z_i=\sqrt{\mu_i/\varepsilon_i}$。切向场连续给出
\[
E_i+E_r=E_t.
\]
对沿 $+z$ 传播的波 $H=E/Z$，沿 $-z$ 传播的反射波磁场方向反转，故
\[
\frac{E_i-E_r}{Z_1}=\frac{E_t}{Z_2}.
\]
定义 $r=E_r/E_i,\ t=E_t/E_i$。由第一式 $t=1+r$，代入第二式：
\[
\frac{1-r}{Z_1}=\frac{1+r}{Z_2},
\]
\[
Z_2-Z_2r=Z_1+Z_1r,
\]
\[
r=\frac{Z_2-Z_1}{Z_2+Z_1},\qquad
t=\frac{2Z_2}{Z_1+Z_2}.
\]
时间平均能流为 $\langle S\rangle=|E|^2/(2Z)$，因此
\[
R=|r|^2,\qquad
T=\frac{Z_1}{Z_2}|t|^2,\qquad R+T=1.
\]
最后一个等式是无吸收条件下的能量守恒检查。
\end{examplebox}

\begin{examplebox}[title=例 6：由推迟势得到简谐偶极子的平均辐射功率]
取 $\bm p(t)=p_0\cos\omega t\,\hat{\bm z}$。远区辐射场为
\[
\bm E_{\rm rad}
=\frac{1}{4\pi\varepsilon_0c^2r}
\left[\hat{\bm r}\times(\hat{\bm r}\times\ddot{\bm p})\right]_{t-r/c},
\qquad
\bm B_{\rm rad}=\frac1c\hat{\bm r}\times\bm E_{\rm rad}.
\]
因 $\abs{\hat{\bm r}\times(\hat{\bm r}\times\hat{\bm z})}=\sin\theta$，
\[
E_{\rm rad}
=\frac{\omega^2p_0}{4\pi\varepsilon_0c^2r}
\sin\theta\cos[\omega(t-r/c)]
\]
（方向沿 $\hat{\bm\theta}$，整体符号不影响功率）。周期平均 Poynting 矢量为
\[
\langle S_r\rangle
=\frac12\varepsilon_0cE_0^2
=\frac{\omega^4p_0^2}{32\pi^2\varepsilon_0c^3r^2}\sin^2\theta.
\]
故单位立体角功率
\[
\dv{\langle P\rangle}{\Omega}
=r^2\langle S_r\rangle
=\frac{\omega^4p_0^2}{32\pi^2\varepsilon_0c^3}\sin^2\theta.
\]
再用
\[
\int\sin^2\theta\,\dd\Omega
=2\pi\int_0^\pi\sin^3\theta\,\dd\theta=\frac{8\pi}{3},
\]
得到
\[
\boxed{\langle P\rangle
=\frac{\omega^4p_0^2}{12\pi\varepsilon_0c^3}}.
\]
$\omega^4$ 依赖说明低频小天线辐射效率很低；$r^{-2}$ 能流保证穿过任意大球面的总功率不变。
\end{examplebox}

\section{量子力学：边界匹配、近似与跃迁}

\begin{examplebox}[title=例 7：有限深方势阱的束缚态条件]
取对称势阱 $V(x)=-V_0$（$|x|<a$），外部 $V=0$。束缚态满足 $-V_0<E<0$，定义
\[
k=\frac{\sqrt{2m(E+V_0)}}{\hbar},\qquad
\kappa=\frac{\sqrt{-2mE}}{\hbar}.
\]
偶宇称态在阱内可写为 $A\cos kx$，右侧为 $Be^{-\kappa x}$。在 $x=a$ 连续：
\[
A\cos ka=Be^{-\kappa a},
\qquad
-Ak\sin ka=-\kappa Be^{-\kappa a}.
\]
第二式除以第一式，消去归一化常数：
\[
\boxed{k\tan ka=\kappa}\qquad\text{（偶态）}.
\]
奇态类似取 $A\sin kx$，得到
\[
\boxed{-k\cot ka=\kappa}\qquad\text{（奇态）}.
\]
同时 $k^2+\kappa^2=2mV_0/\hbar^2$。因此能级由超越方程的交点决定，而不是无限深势阱的简单整数平方。$V_0\to\infty$ 时 $\kappa\to\infty$，边界迫使阱壁处波函数趋零，恢复无限深势阱。
\end{examplebox}

\begin{examplebox}[title=例 8：二能级系统的简并微扰]
未扰 Hamiltonian 在 $\{\ket1,\ket2\}$ 中有共同能量 $E_0$，微扰为
\[
W=\begin{pmatrix}\Delta&v\\v^*&-\Delta\end{pmatrix}.
\]
若直接套用非简并公式，会出现分母 $E_0-E_0=0$。正确做法是在简并子空间中先求
\[
\det(W-\lambda I)
=(\Delta-\lambda)(-\Delta-\lambda)-|v|^2=0,
\]
即
\[
\lambda_\pm=\pm\sqrt{\Delta^2+|v|^2}.
\]
一阶修正后的能量为
\[
E_\pm=E_0\pm\sqrt{\Delta^2+|v|^2}.
\]
当 $\Delta=0$ 时，任意非零耦合 $v$ 都把简并劈裂为 $2|v|$。近自由电子在布里渊区边界的能隙正是同一个二能级结构，其中 $v$ 对应周期势 Fourier 分量 $V_G$。
\end{examplebox}

\begin{examplebox}[title=例 9：方势垒的 WKB 指数与精确结果比较]
高度 $V_0$、宽度 $a$ 的势垒中 $E<V_0$。禁阻区经典动量绝对值
\[
\abs{p(x)}=\sqrt{2m(V_0-E)}\equiv\hbar\kappa .
\]
WKB 透射率的主指数为
\[
T_{\rm WKB}\sim
\exp\left[-\frac2\hbar\int_0^a\abs{p(x)}\,\dd x\right]
=e^{-2\kappa a}.
\]
而边界精确匹配给出
\[
T_{\rm exact}
=\left[
1+\frac{V_0^2\sinh^2(\kappa a)}{4E(V_0-E)}
\right]^{-1}.
\]
当 $\kappa a\gg1$ 时，$\sinh^2(\kappa a)\simeq e^{2\kappa a}/4$，故
\[
T_{\rm exact}\simeq
\frac{16E(V_0-E)}{V_0^2}e^{-2\kappa a}.
\]
WKB 正确抓住指数敏感性，但一般不保证前因子正确。若转折点变化过于尖锐，标准平滑转折点连接公式也需谨慎使用。
\end{examplebox}

\section{热力学与统计物理：从配分函数到响应}

\begin{examplebox}[title=例 10：两能级顺磁体的配分函数、磁化与热容]
每个自旋在磁场中的能量为 $\epsilon_\pm=\mp\mu B$，单粒子配分函数
\[
z=e^{\beta\mu B}+e^{-\beta\mu B}=2\cosh(\beta\mu B).
\]
$N$ 个独立定域自旋有 $Z=z^N$。自由能
\[
F=-Nk_BT\ln[2\cosh(\beta\mu B)].
\]
磁化为
\[
M=-\left(\pdv FB\right)_T
=N\mu\tanh(\beta\mu B).
\]
弱场 $\beta\mu B\ll1$ 时 $\tanh x\simeq x$，于是
\[
\chi=\left(\pdv MB\right)_{B=0}
=\frac{N\mu^2}{k_BT},
\]
即 Curie 定律。内能
\[
U=-\pdv{\ln Z}{\beta}
=-N\mu B\tanh(\beta\mu B),
\]
故
\[
C_B=\left(\pdv UT\right)_B
=Nk_B(\beta\mu B)^2\operatorname{sech}^2(\beta\mu B).
\]
它在高、低温都趋零，中间出现 Schottky 峰；这是一条可直接与比热实验比较的预测。
\end{examplebox}

\begin{examplebox}[title=例 11：一维 Ising 链的转移矩阵]
零外场一维 Ising 链
\[
H=-J\sum_{i=1}^Ns_is_{i+1},\qquad s_i=\pm1,\quad s_{N+1}=s_1.
\]
把每条键的 Boltzmann 因子写成矩阵元
\[
T_{ss'}=e^{\beta Jss'}
=\begin{pmatrix}e^{\beta J}&e^{-\beta J}\\
e^{-\beta J}&e^{\beta J}\end{pmatrix}.
\]
周期边界下
\[
Z=\sum_{\{s_i\}}T_{s_1s_2}\cdots T_{s_Ns_1}
=\operatorname{Tr}T^N.
\]
转移矩阵本征值为
\[
\lambda_+=2\cosh\beta J,\qquad
\lambda_-=2\sinh\beta J,
\]
所以
\[
Z=\lambda_+^N+\lambda_-^N.
\]
热力学极限中 $|\lambda_-|<\lambda_+$，每格自由能为
\[
f=-k_BT\lim_{N\to\infty}\frac1N\ln Z
=-k_BT\ln(2\cosh\beta J).
\]
对任何 $T>0$，该式解析，因此零场一维短程 Ising 模型没有有限温相变。平均场理论若预言有限 $T_c$，其失败来自忽略一维中的强涨落和畴壁熵。
\end{examplebox}

\begin{examplebox}[title=例 12：由能量涨落推出正则系综热容]
正则系综
\[
Z=\sum_i e^{-\beta E_i},\qquad
\langle E\rangle=-\pdv{\ln Z}{\beta}.
\]
再求一次导数：
\[
\pdv{\langle E\rangle}{\beta}
=-\frac{\partial^2\ln Z}{\partial\beta^2}
=-\left[
\frac1Z\frac{\partial^2 Z}{\partial\beta^2}
-\frac1{Z^2}\left(\pdv Z\beta\right)^2
\right].
\]
而
\[
\frac1Z\frac{\partial^2 Z}{\partial\beta^2}=\langle E^2\rangle,\qquad
\frac1Z\pdv Z\beta=-\langle E\rangle,
\]
故
\[
-\pdv{\langle E\rangle}{\beta}
=\langle E^2\rangle-\langle E\rangle^2.
\]
利用 $\dd\beta/\dd T=-1/(k_BT^2)$，
\[
\boxed{\operatorname{Var}(E)
=k_BT^2C_V}.
\]
右边非负立即给出正则系综的 $C_V\ge0$。相对涨落通常随粒子数按 $N^{-1/2}$ 衰减，这解释了宏观热力学量为何看起来确定。
\end{examplebox}

\section{固体物理：结构、能带与输运}

\begin{examplebox}[title=例 13：fcc 单原子结构因子与允许反射]
把常规立方胞中的四个 fcc 格点写为
\[
\bm r_1=(0,0,0),\quad
\bm r_2=\frac a2(0,1,1),\quad
\bm r_3=\frac a2(1,0,1),\quad
\bm r_4=\frac a2(1,1,0).
\]
对 $\bm G=(2\pi/a)(h,k,l)$，若原子散射因子均为 $f$，
\[
S_{hkl}=f\left[
1+e^{-i\pi(k+l)}+e^{-i\pi(h+l)}+e^{-i\pi(h+k)}
\right].
\]
若 $h,k,l$ 全偶，所有指数均为 $1$，$S=4f$；若全奇，两两之和仍为偶数，也有 $S=4f$。若奇偶混合，四项中两正两负，$S=0$。因此 fcc 允许反射条件为
\[
\boxed{h,k,l\ \text{全奇或全偶}}.
\]
消光不是“不满足 Bragg 条件”，而是不同基元散射振幅在满足动量条件后仍发生相消。
\end{examplebox}

\begin{examplebox}[title=例 14：一维紧束缚链的能带和有效质量]
每格一个局域轨道 $\ket n$，只保留最近邻跃迁：
\[
H=\epsilon_0\sum_n\ket n\bra n
-t\sum_n(\ket n\bra{n+1}+\ket{n+1}\bra n).
\]
取 Bloch 叠加
\[
\ket k=\frac1{\sqrt N}\sum_ne^{ikna}\ket n.
\]
作用 Hamiltonian：
\[
H\ket k
=\frac1{\sqrt N}\sum_n
\left[\epsilon_0e^{ikna}-t e^{ik(n-1)a}-t e^{ik(n+1)a}\right]\ket n,
\]
提取 $e^{ikna}$ 得
\[
E(k)=\epsilon_0-t(e^{ika}+e^{-ika})
=\epsilon_0-2t\cos ka.
\]
带底 $k=0$ 附近
\[
E(k)\simeq\epsilon_0-2t+ta^2k^2,
\]
与 $\hbar^2k^2/(2m^*)$ 比较：
\[
m^*=\frac{\hbar^2}{2ta^2}>0.
\]
带顶 $k=\pi/a+q$ 附近曲率为负，电子有效质量为负；改用带正电的空穴描述即可恢复正质量的准粒子图像。
\end{examplebox}

\begin{examplebox}[title=例 15：抛物能带的 Hall 系数]
稳态 Drude 方程在 $\bm B=B\hat{\bm z}$ 中为
\[
0=-e(\bm E+\bm v\times\bm B)-\frac m\tau\bm v.
\]
分量写成
\[
\frac m\tau v_x+eBv_y=-eE_x,\qquad
-eBv_x+\frac m\tau v_y=-eE_y.
\]
Hall 测量中横向开路，$j_y=-nev_y=0$，故 $v_y=0$。第二式给出
\[
E_y=-Bv_x.
\]
又 $j_x=-nev_x$，因此
\[
R_H\equiv\frac{E_y}{j_xB}
=\frac{-Bv_x}{(-nev_x)B}
=-\frac1{ne}.
\]
符号揭示载流子电荷；但在多带、各向异性费米面或不同弛豫时间并存时，$R_H$ 不再简单等于 $-1/(ne)$。回答面试题时必须同时说出这一适用边界。
\end{examplebox}

\section{把推导压缩成面试表达}

\begin{interviewbox}
\textbf{30 秒版的结构：}先说模型与结论，再说决定性一步，最后说适用条件。例如偶极辐射：“局域、非相对论源在远区只保留 $1/r$ 项；由推迟势和电荷守恒把电流积分写成 $\dot{\bm p}$，得到 $\bm E_{\rm rad}\propto\hat{\bm r}\times(\hat{\bm r}\times\ddot{\bm p})/r$，因此角分布为 $\sin^2\theta$，总功率正比于 $\ddot p^{\,2}$。”

\textbf{2 分钟版的结构：}在 30 秒版之间补两行数学：起点、近似、消去中间变量、检查。不要连续背十个结论。

\textbf{黑板版的结构：}左上角写假设和符号；中间只写不可省略的三至五步；右下角框结论并列出一个极限检查。板书空间也是推导逻辑的一部分。
\end{interviewbox}

\begin{selftestbox}
\begin{enumerate}
  \item 若例 1 中支点受到水平简谐势 $KX^2/2$，哪个守恒量消失？重新写两条运动方程。
  \item 用唯一性定理说明镜像法为何只需验证边界，而不必直接计算导体表面每一点的感应电荷。
  \item 从例 6 的角分布解释偶极轴方向为何没有辐射。
  \item 画出例 7 两个超越方程的图像，判断势阱至少存在几个束缚态。
  \item 把例 8 的结果用于 $k=G/2$ 的近自由电子，指出 $\Delta$ 与 $v$ 的物理来源。
  \item 由例 10 求高温极限内能，并解释为何趋于零。
  \item 由例 11 求最近邻关联 $\langle s_is_{i+1}\rangle$。
  \item 给例 14 加上次近邻跃迁 $-t'$，推导新的 $E(k)$。
  \item 若载流子电荷为 $+e$，重做例 15 并判断 Hall 电压符号。
  \item 从十五道例题中任选五道，各写一版 30 秒口述和一版不超过十行的黑板推导。
\end{enumerate}
\end{selftestbox}
